\section{DDP：用局部二次模型改进整条轨迹}
\label{sec:16-Control-and-Reinforcement-Learning/01-Optimal-Control/06}

自动泊车要把车从路边斜着倒进车位，五秒钟，五十个控制步，每步两个量：方向盘角度和油门。前面那套 HJB 的办法在这里用不上——车辆模型带三角函数，值函数没有二次型的形状，状态空间画不成一张表。于是有人提议把五十步的控制拼成一个 100 维向量，代价对它求梯度，直接做梯度下降。梯度确实算得出来，跑起来却极慢，而且经常在快贴上车位时拧成一团。问题出在第一步：那一下打方向的后果，要经过之后四十九步的动力学才显现，梯度把这条长链压缩成一个数字，时序结构在压缩中丢了。

\subsection{把整段动作当成一个长向量，就丢掉了时序结构}

先把问题写下来。离散动力学与有限时域代价：

\[
x_{t+1} = f_t(x_t, u_t),
\qquad
J = \phi(x_T) + \sum_{t=0}^{T-1} \ell_t(x_t, u_t).
\]

朴素的看法把 \(u_{0:T-1}\) 当成一个无结构的决策向量。这个看法没有错，只是浪费。代价对 \(u_0\) 的依赖不是随意的：它先改变 \(x_1\)，再由 \(x_1\) 改变 \(x_2\)，一路传下去。这条链是已知的，写在 \(f_t\) 里。一阶方法看不见它，只看见链末端汇总出来的那个梯度分量。

前面 HJB 那一节的动态规划恰好用的是这条链，只是它要在整个状态空间上把值函数求出来。中间路线由此浮现：不画整张地图，只沿手上这条能跑通的轨迹，在它周围一小圈里做动态规划。差分动态规划（DDP）走的正是这条路。

\subsection{沿一条名义轨迹做二阶展开}

取一条名义轨迹 \((\bar x_t, \bar u_t)\)。它不必优秀，只需可执行——把 \(\bar u\) 喂进真实动力学能跑出 \(\bar x\)。倒车的例子里，一条歪歪扭扭但没撞墙的轨迹就够当起点。

定义局部状态---动作价值

\[
Q_t(x,u) = \ell_t(x,u) + V_{t+1}\bigl(f_t(x,u)\bigr),
\]

其中 \(V_{t+1}\) 是从 \(t+1\) 起算的最优剩余代价的局部近似，末端由 \(V_T = \phi\) 起步。围绕 \((\bar x_t,\bar u_t)\) 把 \(Q_t\) 展到二阶，记偏移 \(\delta x, \delta u\)：

\[
Q_t \approx Q_0 + Q_x^{\trans}\delta x + Q_u^{\trans}\delta u
+ \frac12
\begin{pmatrix}\delta x\\ \delta u\end{pmatrix}^{\trans}
\begin{pmatrix} Q_{xx} & Q_{xu}\\ Q_{ux} & Q_{uu}\end{pmatrix}
\begin{pmatrix}\delta x\\ \delta u\end{pmatrix}.
\]

各块由链式法则从下一时刻的 \(V_{t+1}\) 拼出来：

\[
\begin{aligned}
Q_x &= \ell_x + f_x^{\trans} V'_x, &
Q_u &= \ell_u + f_u^{\trans} V'_x, \\
Q_{xx} &= \ell_{xx} + f_x^{\trans} V'_{xx} f_x + V'_x \cdot f_{xx}, &
Q_{uu} &= \ell_{uu} + f_u^{\trans} V'_{xx} f_u + V'_x \cdot f_{uu},
\end{aligned}
\]

\(Q_{ux}\) 同理。撇号表示 \(t+1\) 时刻的量，最后那些带 \(f_{xx}\)、\(f_{uu}\) 的项是动力学的二阶导数与向量 \(V'_x\) 的张量缩并。

\(Q_{uu}\succ 0\) 时，对 \(\delta u\) 最小化这个二次式有闭式解：

\[
\delta u_t = k_t + K_t\,\delta x_t,
\qquad
k_t = -Q_{uu}^{-1}Q_u,
\qquad
K_t = -Q_{uu}^{-1}Q_{ux}.
\]

回代得到 \(V_t\) 的新二次近似，

\[
V_x = Q_x - K_t^{\trans}Q_{uu}k_t,
\qquad
V_{xx} = Q_{xx} - K_t^{\trans}Q_{uu}K_t,
\]

再往前推一个时刻。整趟反向递推的开销随 \(T\) 线性增长。

两块输出各有身份。\(k_t\) 是对名义动作的开环修正量，与状态无关。\(K_t\) 是局部反馈增益：真实轨迹漂离名义值多少，控制就按 \(K_t\) 补偿多少。第一节的判据在这里派上用场——一个规划器如果只吐出 \(k_t\)，交付的是一条计划；连 \(K_t\) 一起吐出来，交付的才是一条规则。DDP 顺手把反馈也算了出来，这是它比纯轨迹优化多出的一份东西。

\subsection{反向算出来的下降只在模型里成立，得正向验一遍}

反向递推承诺的下降属于那个二次模型，不属于真实系统。所以每轮反向之后要正向跑一遍真实动力学：

\[
x_0^{\text{new}} = x_0, \qquad
u_t^{\text{new}} = \bar u_t + \alpha k_t + K_t\bigl(x_t^{\text{new}} - \bar x_t\bigr),
\qquad
x_{t+1}^{\text{new}} = f_t\bigl(x_t^{\text{new}}, u_t^{\text{new}}\bigr).
\]

步长 \(\alpha\in(0,1]\) 只缩放开环项 \(k_t\)，反馈项 \(K_t\) 不缩——\(\alpha\) 变小意味着新轨迹退回名义轨迹附近，而在名义轨迹附近仍然需要完整的反馈来跟踪。

跑完拿实际代价降幅和模型预测的降幅比。二次模型预测的降幅是

\[
\Delta J(\alpha) = -\alpha\sum_t Q_u^{\trans}Q_{uu}^{-1}Q_u
+ \frac{\alpha^2}{2}\sum_t Q_u^{\trans}Q_{uu}^{-1}Q_u ,
\]

两者之比落在合理区间就接受这条新轨迹，作为下一轮的名义轨迹；差得离谱就减小 \(\alpha\) 重来。若 \(Q_{uu}\) 不正定或病态，往对角上加 \(\mu I\) 再求逆，\(\mu\) 大时步长自动收缩成沿负梯度的一小步。

\begin{center}
\begin{tikzpicture}[>=stealth, every node/.style={font=\small}]
  \draw[->,gray!70] (-0.3,0) -- (8.0,0) node[right,black] {\(x_1\)};
  \draw[->,gray!70] (0,-0.3) -- (0,3.7) node[above,black] {\(x_2\)};

  \foreach \p in {(0.6,0.7),(1.7,1.5),(3.0,2.1),(4.4,2.3),(5.7,1.9),(6.7,1.2)}
    \draw[dashed,gray!60] \p circle (0.5);

  \draw[very thick,blue!55!black] plot[smooth] coordinates
    {(0.6,0.7) (1.7,1.5) (3.0,2.1) (4.4,2.3) (5.7,1.9) (6.7,1.2)};
  \draw[thick,orange!80!black] plot[smooth] coordinates
    {(0.6,0.7) (1.9,1.15) (3.2,1.6) (4.6,1.85) (5.9,1.6) (7.0,0.95)};

  \fill (0.6,0.7) circle (2pt);
  \node[anchor=north east] at (0.65,0.65) {\(x_0\)};
  \fill[red!70!black] (7.3,0.6) circle (2.2pt);
  \node[anchor=west] at (7.4,0.6) {目标};

  \node[blue!55!black,anchor=south] at (3.4,3.05) {名义轨迹};
  \node[orange!70!black,anchor=north] at (4.6,1.3) {改进后的轨迹};

  \draw[->,thick,gray!70!black] (7.0,-0.55) -- (0.9,-0.55);
  \node[anchor=north] at (4.0,-0.6) {反向递推：从终点往回算 \(k_t, K_t\)};
  \draw[->,thick,gray!70!black] (0.9,-1.35) -- (7.0,-1.35);
  \node[anchor=north] at (4.0,-1.4) {正向 rollout：用真实动力学检验};
\end{tikzpicture}

\emph{虚线圈是二次模型可信的那一小圈；算法只在这条管子里做动态规划，不去碰管子外面的状态空间。}
\end{center}

图里那串虚线圈是整个方法的收费点与省钱点。省的是状态空间：管子外面一个格点都不用算。付的是有效范围：新轨迹一旦被推出管子，二次模型的预测就不再可信，线搜索存在的意义正是把每一步按回管子里。

\subsection{它是轨迹空间上的 Newton 法}

\(Q_{uu}\) 是代价对该时刻控制的曲率，\(K_t\) 用曲率把梯度方向拧到正确的朝向，\(\mu I\) 是曲率不可靠时的兜底——这三件事和 \hyperref[sec:09-Convex-Optimization/06-Optimization-Algorithms/02]{``Newton 法用曲率校正方向''} 里的做法逐条对应。DDP 就是 Newton 法搬到轨迹空间的版本，只不过它没有把 Hessian 当成一整块稠密矩阵。

省下的量可以算清楚。把 \(u_{0:T-1}\) 当无结构长向量时，Hessian 是 \(Tm\times Tm\)，分解一次要 \(O(T^3m^3)\)。而代价的时序结构让这个 Hessian 呈块三对角形态，反向递推正是在利用这一点：每步只分解一个 \(m\times m\) 的 \(Q_{uu}\)，总开销 \(O(Tm^3)\)。\(T=50\) 时两者差三个数量级。

iLQR 是 DDP 的常用简化：丢掉上面那两个张量缩并项 \(V'_x\cdot f_{xx}\) 与 \(V'_x\cdot f_{uu}\)，只保留代价的二阶项和线性化的动力学。丢掉之后二阶收敛降为超线性，但省去了动力学的二阶导数，实现代价低得多，工程上大多数时候用的是它。

\subsection{局部改进就是局部改进}

DDP 交付的保证只有一条：代价单调下降到某个驻点。不同的初始轨迹会收敛到不同的解，倒车的例子里向左打方向和向右打方向可能各自收敛，代价相差很远。初始化因此是算法的一部分，不是可以随便填的参数。

约束是另一处麻烦。方向盘打不过极限、油门有饱和，这些盒约束让逐步的最小化变成带约束的二次规划，闭式解消失，需要投影或积极集之类的变体。障碍与接触更糟：碰撞使动力学不可微，Hessian 在那里没有定义，加正则化也补不出来，一般要换成软接触模型或专门的互补约束形式。

维数上也没有免费午餐。反向递推对 \(T\) 是线性的，对状态维数却带着 \(f_x\)、\(V_{xx}\) 这些 \(n\times n\) 矩阵的乘法，几十维还行，几百维就要在导数计算和矩阵操作上仔细做文章。

把这些边界摆好之后，还有一个前提一直没被审视：上面每一个公式里都出现 \(f_t\)，而它被默认为已知。车辆的轮胎侧偏刚度、地面摩擦系数、执行器的延迟，这些量没有哪一个是从说明书上抄来就准的。下一节处理它们的来源——从输入输出数据里把动力学估出来，以及估出来的东西拿去做控制时会发生什么。
